\section{Potenziali termodinamici - ensemble}

\subsection{Equazione di Eulero}

\begin{frame}{Equazione di Eulero $U=-PV+TS+\sum_i\mu_iN_i$}
    Internal energy $U(V,S,N_i)$ is extensive, volume, entropy and $N_i$ are also extensive, $P, T, \mu_i$ are intensive:
\begin{align*}
    &-P=\PDy{V}{U}|_{S,N_i},\ T=\PDy{S}{U}|_{V,N_i},\ \mu_i=\PDy{N_i}{U}|_{S,V}\tag{U homogeneous function}\\
            &\rightarrow U=-PV+TS+\sum_i\mu_iN_i\\
\end{align*}
\end{frame}

\subsection{Ensemble: $\mu$C $U(V,S,N_i)$, C $F(V,T,N_i)$, GC $\Omega(T,V,\mu_i)$.}

\begin{frame}{Potenziali termodinamici}
    Un potenziale termodinamico \'e una funzione di stato ottenuta, eventualmente tramite trasformazioni di Legendre, scegliendo come variabili naturali quelle controllate dall'ensemble; all'equilibrio, sotto i corrispondenti vincoli, il potenziale appropriato \'e stazionario/minimo.
    \begin{columns}[T]
        \begin{column}{0.5\textwidth}
    Ensemble microcanonico: sistema isolato per cui sono fissate le quantit\'a estensive conservate, tipicamente $(N,V,E)$: thermodynamics potential is U
        \begin{align*}
            &dU=-P\,dV+TdS+\sum_i\mu_i\,dN_i\tag{internal energy: $U(V,S,N_i)$}\\
            &-P=\PDy{V}{U}|_{S,N_i}\tag{fissa P}\\
            &T=\PDy{S}{U}|_{V,N_i}\\
            &\mu_i=\PDy{N_i}{U}|_{S,V}
        \end{align*}
    Ensemble Canonico: variabili fissate $V, T, N_i$ il potenziale termodinamico \'e F
        \begin{align*}
            &F(V,T,N_i)=U-TS:\\
            &dF=-P\,dV-S\,dT+\sum_i\mu_i\,dN_i\\
            &-P=\PDy{V}{F}|_{T,N_i}\tag{fissa P}\\
            &-S=\PDy{T}{F}|_{V,N_i}\tag{fissa S}\\
            &\mu_i=\TDy{N_i}{F}|_{V,T,N_j}\tag{fissa $\mu_i$}
        \end{align*}
        \end{column}
        \begin{column}{0.5\textwidth}
    Ensemble Gran-Canonico: variabili fissate $T,V,\mu_i$
        \begin{align*}
            &\Omega(T,V,\mu_i)=F-\sum_i\mu_iN_i=U-TS-\sum_i\mu_iN_i\\
            &\Omega=-PV\tag{by Euler eq}\\
            &d\Omega=-P\,dV-S\,dT-\sum_iN_i\,d\mu_i\\
            &-P=\PDy{V}{\Omega}|_{T,\mu_i}\tag{fissa P}\\
            &-S=\PDy{T}{\Omega}|_{T,\mu_i}\tag{fissa S}\\
            &-N_i=\TDy{\mu_i}{\Omega}|_{V,T,N_j}\tag{fissa $\mu_i$}
        \end{align*}

        \end{column}
    \end{columns}
    
\end{frame}

\subsection{Limite termodinamico: $V\to\infty$}

\begin{frame}{Bulk matter: $\mu C$}    
Nel $\lim_{V\to\infty}$ le grandezze termodinamiche non dipendono da V: bulk matter.
\begin{align*}
    &\epsilon=\frac{U}{V},\ s=\frac{S}{V},\ n_i=\frac{N_i}{V}\\
    &\epsilon(s,n_i)=-P+Ts+\sum_i\mu_in_i\tag{Euler eq}\\
    &T=\PDy{s}{\epsilon}|_{n_i}\\
    &\mu_i=\PDy{\epsilon}{n_i}|_{s,n_j}\\
    &P=-\epsilon+Ts+\sum_i\mu_in_i\tag{Inversione Eulero}\\
    &nn_i\PDof{n_i}(\frac{\epsilon}{n})=n_i\PDy{n_i}{\epsilon}-\epsilon \frac{n_i}{n}=n_i\mu_i-\epsilon \frac{n_i}{n}\tag{trick}\\
    &\Rightarrow \sum_i\mu_in_i=\sum_inn_i\PDof{n_i}(\frac{\epsilon}{n})+\epsilon\\
    &\Rightarrow P=Ts+\sum_inn_i\PDof{n_i}(\frac{\epsilon}{n})\tag{N averall conserved such baryon}
\end{align*}
\end{frame}

\begin{frame}{Bulk matter: $C$}
    \begin{align*}
        &F(T,V,N_i)=U-TS=-PV+\sum_i\mu_iN_i\\
        &\Rightarrow f(T,n_i)=\lim_{V\to\infty}\frac{F}{V}=-P+\sum_i\mu_in_i\\
        &s=-\PDof{T}{f}|_{n_i},\ \mu_i=\PDy{n_i}{f}|_{T,n_j}\\
        &\Rightarrow P=-f+\sum_in_i\PDy{n_i}{f}
    \end{align*}
\end{frame}

\begin{frame}{Bulk matter: $GC$ - Important for hot dense mater}
    \begin{align*}
        &P(T,\mu_i)=-\frac{\Omega}{V}:\ s=\PDy{T}{P}|_{\mu_i},\ n_i=\PDy{\mu_i}{P}|_{T,\mu_j}\\
        &\epsilon(T,\mu_i)=-P+Ts+\sum_i\mu_in_i=-P+T\PDy{T}{P}+\sum_i\mu_i\PDy{\mu_i}{P}
    \end{align*}
\end{frame}

\section{Equilibrio chimico - $\beta$-equilibrio}

\subsection{Conserved quantities}

\begin{frame}{Chemical equilibrium}
    \begin{align*}
        &dU=-P\,dV+TdS+\sum_i\mu_i\,dN_i\\
        &\Rightarrow dU=\sum_i\mu_i\,dN_i\tag{Any work $PdV=0$, any heat exchange $\delta Q=TdS=0$}\\
        &dU=0\Rightarrow\sum_i\mu_i\,dN_i=0\tag{equilibrium}\\
        &\Pproton+\Pelectron\to\Pneutron+\APnue\tag{neutronization}\\
        &\Rightarrow \mu_{\Pnue}+\mu_{\Pneutron}-\mu_{\Pelectron}+\mu_{\Pproton}=0\\
        &\Pneutron\to\Pproton+\Pelectron+\APnue\tag{nucleons+electrons}\\
        &\Rightarrow\mu_{\Pneutron}=\mu_{\Pproton}+\mu_{\Pelectron}+\mu_{\APnue}\\
        &\Rightarrow\mu_{\APnue}=-\mu_{\Pnue},\ \mu_{\APelectron}=-\mu_{\Pelectron}, \mu_{\Pphoton}=0
    \end{align*}
    Chemical potential of anti-particle is negative of that of particle. As P/AP annihilate in pure energy in form of photons $\mu_{Pphoton}=0$.
\end{frame}

\begin{frame}{Number of indep quantities required to describe mixture of baryons and leptons?}
    \begin{columns}[T]
        \begin{column}{0.5\textwidth}
            \begin{block}{Conserved Quantum Numbers}
                \begin{itemize}
                \item Electrodynamics conserve charges
                \item Weak interactions conserve barion B, and lepton L number (Flavor mixing between 3 different neutrino flavors break individual lepton flavor but not total lepton number)
                \item Strong interactions conserve baryon number B and flavor number of Quarks (Up,Down, Strange, Charm, Bottom, Top - Quark number)
            \end{itemize}
            Tutte le reazioni tra barioni e leptoni conservano il numero barionico B, il numero leptonico L e la carica Q: scegliendo i numeri corrispondenti $N_B$, $N_L$, $N_Q$ il sistema quantistico \'e completamente specificato - one can decompose chemical potential of particle $i$ in terms of the quantum numbers of the particle $\mu_i=B_i\mu_B+L_i\mu_L+Q_i\mu_Q$: all we need are the quantum number $N_B, N_L, N_Q$ to define chemical composition or chemical potential (in ensemble gran canonico $\Omega(T,V,\mu_i)$).
            \end{block}
        \end{column}
        \begin{column}{0.5\textwidth}
            \begin{block}{Connection between chemical potential and composition}
                We are in gran-canonical ens.: $\Omega(T,V,\mu_i)$ and $P(T,\mu_B,\mu_L,\mu_Q)=-\frac{\Omega}{V}$; there are additional astrophysical constraints: charge density $n_Q=\frac{N_Q}{V}=0$ - questo fissa il potenziale chimico di carica
                \begin{align*}
                    &n_Q=\PDy{\mu_Q}{P}=\sum_i\PDy{\mu_i}{P}\PDy{\mu_Q}{\mu_i}=\sum_in_iQ_i=0\\
                    &\Rightarrow\PDy{\mu_Q}{P(T,\mu_B,\mu_L,\mu_Q)}=0
                \end{align*}
                Connection between chemical potential and averall charge density is not trivial and depends also on interactions between particles, as all terms in thermodynamic potential that depends on $\mu_Q$ have to be considered.
                For ''cold'' matter neutrinos escape: $\mu_{\Pnue}=\mu_{\Pnum}=\mu_{\Pnut}$, lepton number is not conserved and $\mu_L=0$ \keyword{$\beta$-equilibrio}. For cold degenerate matter $K_BT\ll E_F$: one can set $T=0$ and system can be described in terms of chemical potential $\mu_B$ or baryon density $n_B$ alone, ie $P(T,\mu_B)\to P(\mu_B)$.
                \keyword{$\mu_L=0$ non implica $\mu_e=0$}: $\mu_e=0*\mu_B+1*\mu_L+(-1)*\mu_Q=-\mu_Q$.
            \end{block}
        \end{column}
    \end{columns}
\end{frame}

\begin{frame}{connection between the chemical potential and the overall charge density is a nontrivial}
    \begin{itemize}
        \item Gas non interagente a $T=0$. $n_i=\frac{k_{F,i}^3}{3\pi^2}$, $\mu_i=\sqrt{k_{F,i}^2+m_i^2}$ quindi $n_i(\mu_i)$ e poich\'e $\mu_i=B_i\mu_B+L_i\mu_L+Q_i\mu_Q$ da cui ho $n_Q(\mu_Q)=\sum_iQ_i\mu_i(B_i\mu_B+L_i\mu_L+Q_i\mu_Q)$ e risolvere $n_Q(\mu_Q)=0$
        \item Relativistic Mean Field theory - Si parte da lagrangiana in cui protoni e neutroni interagiscono tramite campi mesonici, $\sigma$ attrattivo, $\omega_{\mu}$ repulsivo, $\rho_{\mu}$ isovettoriale. In approssimazione di campo medio:
            \begin{align*}
                &\Omega=\Omega_{Fermions}(T,\mu_i^*,m_i^*)+\Omega_{fields}(\sigma,\omega_0,\rho_0)\\
                &\PDy{\sigma}{\Omega}=0,\ \PDy{\omega_0}{\Omega}=0,\ \PDy{\rho_0}{\Omega}=0\tag{stazionariet\'a potenziale determina campi}
            \end{align*}
    \end{itemize}
\end{frame}

\subsection{$\beta$-equilibrium}

\begin{frame}{Gas \Pproton/\Pelectron: sotto $n_c$ domina $m_p+m_e<m_n$. }
    \begin{itemize}
        \item Neutronization - The weak reaction $\Pelectron+\Pproton\to\Pneutron+\Pnue$ proceed only if $\Delta m=m_n-m_p=(393.57-928.28)\si{\mega\ev}=\SI{1.29}{\mega\ev}$ is compensated by kinetic energy.
        \item $\beta$ decay - $\Pneutron\to\Pproton+\Pelectron+\APnue$ is blocked if density is high enough such energy level are occupied.
    If we assume chemical equilibrium, with $T=0$, $\mu_{\nu}=0$: $\mu_e+\mu_p=\mu_n$. A $T=0$ tutti i livelli sono occupati fino all'energia di Fermi: $k_{F,i}$ impulso di Fermi, $g_i=2s_i+1$
    \begin{align*}
        &\mu_i=E_{F,i}=\sqrt{k_{F,i}^2+m_i^2}\tag{$T=0$: energy needed to take fermion}\\
        &\sqrt{k_{F,e}^2+m_e^2}+\sqrt{k_{F,p}^2+m_p^2}=\sqrt{k_{F,n}^2+m_n^2}\tag{equilibrium}\\
        &n_i=g_i\int_0^{k_{F,i}}\frac{d^3k}{(2\pi)^3}=\frac{g_i}{6\pi^2}k_{F,i}^3\tag{number density}\\
        &n_e=n_p\Rightarrow g_ek_{F,e}^3=g_pk_{F,p}^3\tag{charge neutrality}\\
        &g_e=g_p=2\Rightarrow k_{F,e}=k_{F,p}\tag{thermod. quant. as f. of $\mu_B$/$n_B$}\\
        &P=P_e+P_p+P_n\\
        &\epsilon=\epsilon_e+\epsilon_p+\epsilon_n\\
        &n_B=n_p+n_n\\
        &n_Q=n_p-n_e=0
    \end{align*}
                \item A gas of \Pproton-\Pelectron is energ-favor. at low density. $m_p+m_e=\SI{938.79}{\mega\ev}<m_n=\SI{939.57}{\mega\ev}$.
    \end{itemize}
\end{frame}

\begin{frame}{Gas neutroni/protoni/elettroni: neutronization. Formazione nuclei}
            \begin{itemize}
                \item At some critical $n_e=n_p$ ($n_Q=n_p-n_e=0$) Fermi energy of electrons and protons sum up to \Pneutron rest mass ($k_{F,n}=0$):
                    \begin{align*}
                        &\sqrt{k_{F,e}^2+m_e^2}+\sqrt{k_{F,p}^2+m_p^2}=m_n\tag{\Pneutron at rest}\\
                        &K_{F,p}\ll m_p\Rightarrow\sqrt{k_{F,p}^2+m_p^2}\approx m_p\tag{\Pproton NR}\\
                        &\Rightarrow\sqrt{k_{F,e}^2+m_e^2}\approx m_n-m_p\tag{$n_e$ neutronization}\\
                        &\Rightarrow n_c=\frac{1}{3\pi^2}((m_n-m_p)^2+m_e^2)\expy{\frac{3}{2}}\approx\SI{7.4e30}{\per\cubic\cm}
                    \end{align*}
    Con aumentare della densit\'a si ha $n_n\gg n_p$ e si forma la materia neutronica.
        \item Condizione asintotica equilibrio $K_F\gg m_N$ (limite UR). La condizione $\sqrt{k_{F,e}^2+m_e^2}+\sqrt{k_{F,p}^2+m_p^2}=m_n$ si riduce a $k_{F,n}\approx k_{F,e}+k_{F,p}=2k_{F,p}$: essendo $n\propto gk_F^3$ con $g_n=g_p=2$ e quindi nel limite di alte densit\'a $n_n=8n_p$ o $\frac{n_n}{n_B}=\frac{8}{9}$.
    \end{itemize}
    La materia densa, oltre $\rho_c=m_pn_c\approx\SI{1.2e7}{\gram\per\cubic\cm}$, in $\beta$-equilibrio \'e composta principalmente da neutroni in quanto la neutralit\'a imporrebbe di aggiungere coppia \Pproton/\Pelectron, sopra la densit\'a critica non \'e energeticamente conveniente.
\end{frame}

\begin{frame}{Realistic neutron matter: formazione nuclei}
    \begin{itemize}
        \item In realt\'a lo stato di minimo dei nucleoni \'e sono i nuclei: materia in $\beta$-equilibrio consiste di nuclei in gas di elettroni che garantisce neutralit\'a, ma i nuclei presenti non sono quelli con minor massa per barione. We need to find nucleus (A,Z) scanning through nuclear chart which minimizes total energy of the system for given baryon number density. Vedi termine $\frac{Z}{A}\mu_e$ dovuto agli elettroni per neutralit\'a di carica.
    \item In addition nuclei will form regular lattice to minimize coulomb energy, ie lattice energy has to be included; lattice energy is important for composition of nuclei not for EOS.
    \item At low density $^{56}Fe$, the one with lowest mass per baryon, is present in $\beta$-equilibrated matter; with increasing baryon density nuclei become more neutron rich (\xaumenta{n_B}, \xaumenta{n_e}, aumenta probabilit\'a di convertire protone in neutrone) until highest possible \Pneutron/\Pproton ration is achieved: there the baryon density is called neutron-drip density $\rho\approx\SI{4e11}{\gram\per\cubic\cm}$.
    \item Neutron drip density depends upon nuclear model used to determine binding energy of unknown neutron rich nucleus (in particular the ones close to neutron dripline)
    \item Nuclei con B streaordinariamente alta saranno favoriti in materia in $\beta$-equilibrio: closed shell of neutrons to a magic number $N=2,8,20,50,82,126$; modelli indicano come pi\'u favoriti $N=50,82$.
    \end{itemize}

\end{frame}

\begin{frame}{Composizione materia in $\beta$-equilibrio}\frameinlbftrue
    \begin{itemize}
        \item $\rho<\SI{e4}{\gram\per\cubic\cm}$: gas of atoms - Atmospheres of compact stars.
        \item $\rho=\SI{e4}{\gram\per\cubic\cm}$: si forma gas di elettroni liberi e reticolo di nuclei di $^{56}Fe$.
        \item $\rho=\SI{2e6}{\gram\per\cubic\cm}$: elettroni diventano relativistici.
        \item $\rho<\SI{8e6}{\gram\per\cubic\cm}$: si forma un reticolo di nuclei oltre $^{56}Fe$ con attorno gas di elettroni; i nuclei diventano sempre pi\'u neutron rich all'aumentare della densit\'a.
        \item $\rho<\SI{4e11}{\gram\per\cubic\cm}$: neutron drip density - densit\'a limite per emissione di neutroni da èparte di neutron rich nuclei.
    \end{itemize}
\end{frame}

\section{EOS fermionic matter at $T=0$}

\begin{frame}{}
    \begin{columns}[T]
        \begin{column}{0.5\textwidth}
            \begin{itemize}
                \item EOS for degenerate matter in compact star: $P(\epsilon)$ or $P(n)$
                \item EOS to describe nuclear matter: binding energy per baryon as a function of baryon number density $\frac{E_b}{A}(n_B)=\frac{\epsilon}{n_B}-m_N$
                \item Energy per particle $\frac{E}{A}$ is a thermodynamic potential in T.L. function of number density and entropy per particle: possiamo usare formula di Eulero $\epsilon=Ts-P+\mu_Bn_B\xrightarrow{T=0}-P+\mu_Bn_B$ e $\mu_B=\frac{\epsilon}{n_B}$ o direttamente $P=Ts+\sum_inn_i\PDof{n_i}(\frac{\epsilon}{n})$.
            \end{itemize}
        \end{column}
        \begin{column}{0.5\textwidth}
            \begin{align*}
                &\PDof{n_B}(\frac{E}{A})=\PDof{n_B}(\frac{\epsilon}{n_B})=\frac{1}{n_B}\frac{\epsilon}{n_B}-\frac{\epsilon}{n_B^2}\\
                &\Rightarrow\PDof{n_B}(\frac{\epsilon}{A})=\frac{\mu_B}{n_N}-\frac{\epsilon}{n_B^2}=\frac{P}{n_B^2}\\
                &\Rightarrow P=n_B^2\PDof{n_B}(\frac{E}{A})
            \end{align*}
            Energy per particle as function of number density is proportional to pressure: a minimum in energy per particle corresponds to vanishing pressure, ie stable equilibrated matter.
        \end{column}
    \end{columns}
    \begin{itemize}
        \item Le espressioni $P(\epsilon)$ o $P(n)$ non sono potenziali termodinamici e non contengono tutte le informazioni termodinamiche al contrario dell'espressione $\frac{E_B}{A}$: infatti calcolando la densit\'a di energia usando EOS nella forma $P(n)$ tramite la precedente equazione per P la costante di integrazione resta indeterminata; questo non \'e il caso per $P(\mu,T)$:
            \begin{align*}
                &\epsilon(\mu,T)=-P(\mu,T)+\mu n+Ts\\
                &n=(\PDy{\mu}{P(\mu,T)}),\ s=(\PDy{T}{P(\mu,T)})
            \end{align*}
            Infatti nell'ensemble grancanonico $P=-\frac{\Omega}{V}$ direttamente legata a potenziale termodinamico (funzione di stato dalle cui derivate ottengo tutte le grandezze termodinamiche).
    \end{itemize}
\end{frame}

\section{Dense Matter - Fermi gas - Polytropes}

\begin{frame}{Fermi Gas of electrons and nucleons: NR limit}
    \begin{columns}[T]
        \begin{column}{0.45\textwidth}
            \begin{align*}
                &P_e=\frac{g_e}{3(2\pi)^3}\int\,d^3k\scap{k}{v}_e\tag{electrons P}\\
                &v_e=\frac{k}{E_e(k)},\ E_e(k)=\sqrt{k^2+m_e^2}\\
                &\Rightarrow P_E=\frac{g_e}{6\pi^2}\int_0^{k_{F,e}}dk \frac{k^4}{E_e(k)}\\
                &\epsilon=\frac{g}{(2\pi)^3}\int d^3kE(k)
            \end{align*}
        \end{column}
        \begin{column}{0.5\textwidth}
            \begin{align*}
                &T^{ij}=g_e\int\frac{d^3k}{(2\pi)^3}k^iv^j f(\vec{k})\\
                &T^{xx}=T^{yy}=T^{zz}=P\\
                &T^{xx}+T^{yy}+T^{zz}=\frac{g_e}{(2\pi)^3}\int\,d^3k\scap{k}{v}_e\\
                &\dot{x}=\PDy{\vec{p}}{H} (E=\hbar\omega, \vec{p}=\hbar\vec{k},\vec{v}_g=\PDy{\vec{k}}{\omega})\Rightarrow v_e=\nabla_{\vec{v}}E_e
            \end{align*}
        \end{column}
    \end{columns}
    \begin{itemize}
        \item Limite NR $k_{F,e}\ll m_e$, $v_e\approx \frac{k}{m_e}$, $E_e\approx m_e$: $P_e^{NR}=\frac{g_e}{6\pi^2}\int_0^{k_{F,e}}dk \frac{k^4}{m_e}=\frac{g_e}{30\pi^2}\frac{k_{F,e}^5}{m_e}$
        \item Number density $n_e=\frac{g_e}{(2\pi)^3}\int d^3k=\frac{g_e}{2\pi^2}\int_0^{k_{F,e}}k^2\,dk=\frac{g_e}{6\pi^2}k_{F,e}^3$
        \item $P_e^{NR}=\frac{1}{5}(\frac{6\pi^2}{g_e})\expy{\frac{2}{3}}\frac{n_e\expy{\frac{5}{3}}}{m_e}$: pressure inversely proportional to fermion mass, the contribution from nucleons is suppressed by a factor $\frac{m_e}{m_N}$ (from nucleons is suppressed by a factor $\frac{m_e}{Am_N}$) and can be ignored.
        \item Charge neutrality: $n_e=n_p=Zn_A=\frac{Z}{A}n_N$.
        \item Energy in NR limit $E(k)\approx m$: $\epsilon^{NR}\approx m \frac{g}{(2\pi)^3}\int d^3k=mn=\rho$, quindi la densit\'a di energia \'e dominata dai nucleoni $\epsilon^{NR}=\rho_A=m_A\rho_A\approx\rho_N=m_Nn_N$ and ignoring binding energy we have $m_A\approx Am_N$.
        \item Total pressure in terms of mass density: $P^{NR}=\frac{1}{5}(\frac{6\pi^2}{g_e})\expy{\frac{2}{3}}\frac{1}{m_e}(\frac{Z\rho_N}{Am_N})\expy{\frac{5}{3}}\propto(\frac{Z}{A})\expy{\frac{5}{3}}\rho_N\expy{\frac{5}{3}}$ - EOS in NR limit is a power law with exponent $\frac{5}{3}$.
    \end{itemize}
\end{frame}

\begin{frame}{Fermi Gas of electrons and nucleons: R limit}
    \begin{columns}[T]
        \begin{column}{0.45\textwidth}
            \begin{align*}
                &P_e=\frac{g_e}{3(2\pi)^3}\int\,d^3k\scap{k}{v}_e\tag{electrons P}\\
                &v_e=\frac{k}{E_e(k)},\ E_e(k)=\sqrt{k^2+m_e^2}\\
                &\Rightarrow P_E=\frac{g_e}{6\pi^2}\int_0^{k_{F,e}}dk \frac{k^4}{E_e(k)}\\
                &\epsilon=\frac{g}{(2\pi)^3}\int d^3kE(k)
            \end{align*}
        \end{column}
        \begin{column}{0.5\textwidth}
            \begin{align*}
                &T^{ij}=g_e\int\frac{d^3k}{(2\pi)^3}k^iv^j f(\vec{k})\\
                &T^{xx}=T^{yy}=T^{zz}=P\\
                &T^{xx}+T^{yy}+T^{zz}=\frac{g_e}{(2\pi)^3}\int\,d^3k\scap{k}{v}_e\\
                &\dot{x}=\PDy{\vec{p}}{H} (E=\hbar\omega, \vec{p}=\hbar\vec{k},\vec{v}_g=\PDy{\vec{k}}{\omega})\Rightarrow v_e=\nabla_{\vec{v}}E_e\\
                &n_e=n_p=Zn_A=\frac{Z}{A}n_N\tag{charge netrality}
            \end{align*}
        \end{column}
    \end{columns}
    \begin{itemize}
        \item Limite R $k_{F,e}\gg m_e$, $k_{F,e}\ll m_N$: relativistic \Pelectron, non-relativistic nucleon. $v_e\approx c=1$, $E^{R}_e\approx k$:             
            \begin{align*}
                &P_e^R=\frac{g_e}{6\pi^2}\int_0^{k_{F,e}}dk k^3=\frac{g_e}{24\pi^2}k_{F,e}^4=\frac{1}{4}(\frac{6\pi^2}{g_e})\expy{\frac{1}{3}}n_e\expy{\frac{4}{3}}\\
                &P_e^R=\frac{g_e}{24\pi^2}k_{F,e}^4=\frac{1}{4}n_ek_{F,e},\ P_N^{NR}=\frac{g_N}{30\pi^2}\frac{k_{F,N}^5}{m_N}=\frac{1}{5}n_N \frac{k_{F,N}^2}{m_N}\Rightarrow \frac{P_p^{NR}}{P_e^R}=\frac{4}{5}\frac{k_{F,e}}{m_N}\ll1\\
                &(n_e=\frac{g_e}{(2\pi)^3}\int d^3k=\frac{g_e}{2\pi^2}\int_0^{k_{F,e}}k^2\,dk=\frac{g_e}{6\pi^2}k_{F,e}^3,\ n_e=n_p,\ g_p=g_e=2)
            \end{align*}
        \item Energy in R limit $E(k)\approx k$: $\epsilon^R=\rho_A=m_A\rho_A\approx\rho_N=m_Nn_N$ and ignoring binding energy we have $m_A\approx Am_N$.
        \item EOS: $P^R_e=\frac{1}{4}(\frac{6\pi^2}{g_e})\expy{\frac{1}{3}}(\frac{Z\rho_N}{Am_N})\expy{\frac{4}{3}}\propto(\frac{Z}{A})\expy{\frac{4}{3}}\rho_N\expy{\frac{4}{3}}$.
    \end{itemize}
\end{frame}

\begin{frame}{EOS pure neutron gas}
    \begin{columns}[T]
        \begin{column}{0.53\textwidth}
            \begin{align*}
                &P_n^{NR}=\frac{g_n}{6\pi^2}\int_0^{k_{F,n}}dk \frac{k^4}{m_n}=\frac{g_n}{30\pi^2}\frac{k_{F,n}^5}{m_n}\tag{NR: $k_F\ll m_n$}\\
                &\epsilon^{NR}_n=\rho_n=m_nn_n\ (n_n=\frac{g_n}{(2\pi)^3}\int d^3k=\frac{g_n}{6\pi^2}k_{F,n}^3)\\
                &P^{NR}_n=\frac{1}{5}(\frac{6\pi^2}{g_n})\expy{\frac{2}{3}}\frac{1}{m_n}(\frac{\rho_n}{m_n})\expy{\frac{5}{3}}\propto\rho_n\expy{\frac{5}{3}}\tag{EOS}
            \end{align*}
        \end{column}
        \begin{column}{0.47\textwidth}
            \begin{align*}
                &P_e^R=\frac{g_n}{6\pi^2}\int_0^{k_{F,n}}dk k^3\tag{UR: $k_{F,n}\gg m_n$}\\
                &=\frac{g_n}{24\pi^2}k_{F,n}^4=\frac{1}{4}(\frac{6\pi^2}{g_n})\expy{\frac{1}{3}}n_n\expy{\frac{4}{3}}\ (n_n=\frac{g_n}{6\pi^2}k_{F,n}^3)\\
                &\epsilon=\frac{g}{(2\pi)^3}\int d^3kE\xrightarrow{E_n^R\approx k}\frac{g_n}{(2\pi)^3}\int kd^3k\tag{$\epsilon^R_n$}\\
                &\Rightarrow\epsilon^R_n=\frac{g_n}{2\pi^2}\int_0^{k_{F,n}}dkk^3=\frac{g_n}{8\pi^2}k_{F,n}^4
            \end{align*}
            EOS approaches in UR limit the EOS of free gas of massless particles $P^R=\frac{1}{3}\epsilon^R$
        \end{column}
    \end{columns}
\end{frame}

\begin{frame}{Parametrization EOS of free fermion gas - Politrope}
    \begin{itemize}
        \item Parametrization of free fermion gas EOS: $P\propto n^{\Gamma}=(\frac{\rho}{m})^{\Gamma}$.
        \item For the TOV we need relation between pressure-energy density: ansaz $P\propto\epsilon^{\gamma}$.
        \item Non Relativistic limit: $P\propto n\expy{\frac{5}{3}}\propto\epsilon\expy{\frac{5}{3}}$ quindi $\Gamma=\gamma=\frac{5}{3}$.
        \item Relativistic limit: $P\propto n\expy{\frac{4}{3}}\propto\epsilon\expy{\frac{4}{3}}$ quindi $\Gamma=\gamma=\frac{4}{3}$.
        \item Ultra-relativistic limit for single fermion gas: $P\propto n\expy{\frac{4}{3}}$, $P=\frac{1}{3}\epsilon$ quindi $\Gamma=\frac{4}{3}$, $\gamma=1$.
    \end{itemize}
    $\rho$ stands for density of non-relativistic species, $\epsilon$ denotes the more general energy density.
\end{frame}

\begin{frame}{EOS of free Fermi gas in analitic parametric form}
    \begin{align*}
        &P=\frac{g}{3(2\pi)^3}\int d^3k \frac{k^2}{E(k)}=\frac{g}{2\pi^2}\int_0^{k_F}dk \frac{k^4}{\sqrt{k^2+m^2}}=\frac{g}{24\pi^2}\left\{k_FE_F(k_F^2-\frac{3}{2}m^2)+\frac{3}{2}m^4\ln{\frac{k_F+E_F}{m}}\right\}\\
        &\epsilon=\frac{g}{(2\pi)^3}\int d^3kE(k)=\frac{g}{2\pi^2}\int_0^{k_F}dkk^2\sqrt{k^2+m^2}=\frac{g}{8\pi^2}\left\{k_FE_F(k_F^2+\frac{1}{2}m^2+\frac{1}{2}m^4\ln{\frac{k_F+E_F}{m}})\right\}
    \end{align*}
    For given Fermi moment $k_F$ or number density n, using $n=\frac{g}{2\pi^2}k_F^3$, pressure and energy density can be computed parametrically.
\end{frame}

\begin{frame}{Politrope}
    \begin{itemize}
        \item $P=Kn^{\Gamma}$, dimension of K depends on $\Gamma$. Extension to density dependent $\Gamma=\PDly{n}{P}$, usefull in transition from NR to R regimes.
        \item We can't derive $\epsilon$ and $\mu$ from polytrope: i potenziali termodinamici sono della forma $\epsilon(n)$ o $P(\mu)$.
        \item Invertendo $P=n^2\TDy{n}{\frac{\epsilon}{n}}$ la densit\'a di energia pu\'o essere riscritta in funzione della pressione $\TDy{n}{\frac{\epsilon}{n}}=\frac{P}{n^2}=Kn^{\Gamma-2}$ e integrando $\frac{\epsilon}{n}=\frac{K}{\Gamma-1}n^{\Gamma-1}+\const{}$, compare una costante da determinare, per limite bassa densit\'a l'energia per particella $\frac{E}{N}$ \'e la massa della particella $\frac{E}{N}=\frac{\epsilon}{n}\to m$ for $n\to0$, quindi la costante di integrazione pu\'o essere associata alla massa della particella per $\Gamma>1$, quindi la densit\'a di energia $\epsilon=mn+\frac{K}{\Gamma-1}n^{\Gamma}$; adesso possiamo calcolare $\mu=\TDy{n}{\epsilon}=m+\frac{K\Gamma}{\Gamma-1}n^{\Gamma-1}$
        \item Acausal behaviour at high density: at high density speed of sound in the medium, for polytropic EOS, can exceed speed of light.
            \begin{align*}
                &c_s^2=\PDy{\epsilon}{P}|_s=\PDy{n}{P}\TDy{\epsilon}{n}\to\Gamma-1\tag{$n\to\infty$}
            \end{align*}
            For $\Gamma=2$, the stiffest EOS compatible with causality, $c_s=1$ and $P=\epsilon$; note that for a fermi gas in UR limit, ie for high density, $c_s=\frac{1}{\sqrt{3}}$
    \end{itemize}
\end{frame}

\section{Superfluidit\'a}

\begin{frame}{Biblio}
    \begin{itemize}
        \item An introduction to Ginzburg-Landau theory of phase-transitions and non-equilibrium patterns - Hohenberg 2015: Review
        \item Introduction to Superconductivity - Tinkham: Ch 1 Historical Overview, Ch 2 Introduction to electrodynamics of superconductors, Ch 3 The BCS theory, Ch 4 Ginzburg-Landau Theory
    \end{itemize}
\end{frame}

\begin{frame}{Relazione tra Landau-Ginzburg e BCS}
    \begin{itemize}
        \item BCS (1957) \'e una teoria microscopica del pairing, Ginzburg-Landau (1950) \'e sua descrizione efficace macroscopica vicino a $T_c$; Gor'kov mostr\'o che le LG possono essere ricavate da BCS.
        \item In BCS l'oggetto microscopico centrale \'e il parametro di pairing $\Delta(\vec[\mathbf] r) \propto \pairfield$
        \item In GL si introduce fenomenologicamente un parametro d'ordine complesso $\psi_{GL}(\vec[\mathbf]{r})=|\psi_{GL}|\exp{i\theta}$: vicino alla temperatura di transizione $\psi_{GL}(\vec[\mathbf] r)\propto\Delta(\vec[\mathbf] r)$
    \end{itemize}
\end{frame}

\subsection{Superconductivity/Superfluidity: Fenomenologia}

\begin{frame}{Fenomenologia superconduttori}
    \begin{itemize}
        \item La resistenza di vari materiali si annulla per $T<T_c$ dipendente dal materiale.
        \item Meissner Effect: campo magnetico \'e espulso da campione per $T<T_c$: questo non \'e spiegato da resistenza nulla che tenderebbe a ''congelare'' il campo magnetico. Superconductivity will be destroied by critical magnetic field $H_c$: $H_c$ is determined equating magnetic energy density to $\frac{H^2}{8\pi}$ to Helmoholtz free-ergy difference between normal/condensed state $\frac{H^2}{8\pi}=f_n(T)-f_s(T)$ where $f$ is Helmholtz free-energy per unit volume in the respective phase with zero-field. Approximated by parabolic law $H_c(T)\approx H_c(0)[1-(\frac{T}{T_c})^2]$.
        \item The transition a zero field at $T_c$ is of second order, the transition in the presence of field is of first order since there is a discontinuos change in thermodynamics state of system and associated latent heat.
        \item London eqs: $n_s$ numeric density of superconducting electrons, $\vec{h},\vec{e}$ microscopic fields
            \begin{columns}[T]
                \begin{column}{0.5\textwidth}
                    \begin{align*}
                        &\vec{E}=\PDy{t}{\Lambda\vec{J}_s},\ \vec{h}=-c\nabla\wedge(\Lambda\vec{J}_s),\ \Lambda=\frac{4\pi^2\lambda^2}{c^2}=\frac{m}{n_se^2}\\
                        &\vec{p}\to m\vec{v}+\frac{e\vec{A}}{c}\Rightarrow\exv{\vec{v}_s}=-\frac{a\vec{A}}{mc},\ \vec{E}=\nabla\phi-\frac{1}{c}\PDy{t}{\vec{A}}\\
                        &\vec{J}_s=n_se\exv{\vec{v}_s}=-\frac{n_se^2\vec{A}}{mc}=-\frac{\vec{A}}{\Lambda c}
                    \end{align*}
                \end{column}
                \begin{column}{0.5\textwidth}
                    $\vec{e}$: Any electric field accelerates superconducting electrons. $\vec{h}$: since $\nabla\wedge\vec{h}=\frac{4\pi\vec{J}}{c}$ we have exponentilly screened interior, $\nabla^2\vec{h}=\frac{\vec{h}}{\lambda^2}$, with penetration depth $\lambda$. In absence of field we have zero net momentum of GS. Time derivative of eq for $\vec{J}_s$ give $\vec{e}$, curl gives $\vec{h}$. $\lambda_L(0)=\sqrt{\frac{mc^2}{4\pi ne^2}}$.
                \end{column}
            \end{columns}
        \item Pippard: coherence length $\xi_0$ of s. wf - $\vec{J}(\vec{r})=\frac{3\sigma}{4\pi l}\int \frac{\vec{R}[\vec{R}\cdot\vec{R}(\vec{r'})]\exp{-\frac{R}{l}}}{R^4}\,d\vec{r'}$, $\vec{R}=\vec{r}-\vec{r'}$ (Chamber's Ohm law gen.)
            \begin{columns}[T]
                \begin{column}{0.5\textwidth}
                    \begin{align*}
                        &\Delta p\approx \frac{kT_c}{v_f}\Rightarrow\Delta x\geq \frac{\hbar}{\Delta p}\approx \frac{\hbar v_F}{kT_c}\Rightarrow \xi_0=a\frac{\hbar v_F}{kT_c}\\
                        &\vec{J}_s(\vec{r})=\frac{3}{4\pi\xi_0\Lambda c}\int \frac{\vec{R}[\vec{R}\cdot\vec{A}(\vec{r'})]\exp{-\frac{R}{\xi}}\,d\vec{r'}}{R^4},\ \frac{1}{\xi}=\frac{1}{\xi_0}+\frac{1}{l}
                    \end{align*}
                \end{column}
                \begin{column}{0.5\textwidth}
                    Current at $\vec{r}$ deps on $\vec{E}(\vec{r'})$ throughout volume radius $\approx l$ about $\vec{r}$ - typic. $\xi_0\gg\lambda_L(0)$ - $\xi_0$ smallest size wave packet that superconductin charge carriers can form; BCS: $a\approx0.18$.
                \end{column}
            \end{columns}
    \end{itemize}
\end{frame}

\begin{frame}{Fenomenologia Superfluidi/superconduttori}
    \begin{block}{Specific Heat}
        Superconductors: energy gap $\Delta\approx kT_c$ between GS and quasi-particle excitations - Specific heat of electrons for $T\ll T_c$ dominated by exponential $C_{es}\approx \gamma T_ca\exp{-b \frac{T_c}{T}}$ where normal state elettronic specific heat is $C_{en}=\gamma T$: exponential deps. ($b\approx1.5$) implies min. excitation energy per particle of $\approx1.5kT_C$ - popolazione eccitazione soppressa di fattore di Boltzman $\exp{-\frac{\Delta}{kT}}$. Il calore specifico misura la dipendenza dell'energia immagazzinata nelle eccitazioni da T
    \end{block}
    \begin{block}{EM absorption in region $\hbar\omega\approx kT_c$}
        Far-IR and microwave measurements shows energy gap $\approx3-4kT_c$: questo risultato \'e consistente con quello calorimetrico solo se le eccitazioni sono prodotte in coppia come atteso se obbescono alla statistica di Fermi.
    \end{block}
    \begin{block}{BCS Theory - Cooper pair with spatial extension $\xi_0$ are superconducting charge carriers}
        Weak attractive interaction between electrons, such second order phonon-electron int., causes instability in Fermi's sea GS to formation of electron bound pairs occuping states with opposite momentum and spin.
        La teoria prevede energia minima $E_g=2\Delta(T)$ per rompere la coppia creando due eccitazioni di quasi-particella; $\Delta(T_c)=0$ e aumenta per $T\ll T_c$ con $E_g(0)=2\Delta(0)=3.528kT_c$
    \end{block}
    \begin{block}{GL Theory - $\psi$ parametro d'ordine}
        \begin{columns}[T]
            \begin{column}{0.5\textwidth}
                \begin{align*}
                    &f_s-f_n=a(T)|\psi|^2+\frac{b}{2}|\psi|^4(+c|\psi|^6+\ldots)\\
                    &
                \end{align*}
            \end{column}
            \begin{column}{0.5\textwidth}
                Local density of superconducting electrons $n_s=|\psi|^2$ - Transizione continua: $\alpha(T-T_c)$; minimo a $|\psi|^2=-\frac{\alpha}{\beta}$
            \end{column}
        \end{columns}
        
    \end{block}
\end{frame}

\begin{frame}{MQ Wizard}
    \begin{block}{Spazio di Hilbert per stato a due particelle a partire da stati a particella singola}
    Per un sistema costituito da due particelle, lo spazio di Hilbert totale e il prodotto tensoriale degli spazi a singola particella $\mathcal{H}=\mathcal{H}_1\otimes\mathcal{H}_2$.

    Se la prima particella si trova nello stato $\ket{\vb{k}_1}\in\mathcal{H}_1$ e la seconda nello stato $\ket{\vb{k}_2}\in\mathcal{H}_2$, il corrispondente stato prodotto a due particelle \'e $\ket{\vb{k}_1,\vb{k}_2}=\ket{\vb{k}_1}\otimes\ket{\vb{k}_2}$. Nella rappresentazione delle coordinate si ha $\Psi(\vb{r}_1,\vb{r}_2)=\braket{\vb{r}_1,\vb{r}_2|\vb{k}_1,\vb{k}_2}$, con $\ket{\vb{r}_1,\vb{r}_2}=\ket{\vb{r}_1}\otimes\ket{\vb{r}_2}$. Poich\'e il prodotto scalare sullo spazio prodotto fattorizza $\braket{\vb{r}_1,\vb{r}_2|\vb{k}_1,\vb{k}_2}=\braket{\vb{r}_1|\vb{k}_1}\braket{\vb{r}_2|\vb{k}_2}$, la funzione d'onda a due particelle \'e quindi il prodotto delle funzioni
d'onda a singola particella $\Psi(\vb{r}_1,\vb{r}_2)=\phi_{\vb{k}_1}(\vb{r}_1)\phi_{\vb{k}_2}(\vb{r}_2)$. Nel caso di stati a onda piana $\phi_{\vb{k}}(\vb{r})=\frac{1}{\sqrt{\Omega}}e^{i\vb{k}\cdot\vb{r}}$, segue $\Psi(\vb{r}_1,\vb{r}_2)=\frac{1}{\Omega}e^{i\vb{k}_1\cdot\vb{r}_1}e^{i\vb{k}_2\cdot\vb{r}_2}$. L'utilizzo di stati prodotto non implica che l'interazione tra le particelle sia nulla: per $V\neq0$ essi possono essere usati come base di espansione, mentre gli autostati dell'Hamiltoniana completa sono in generale sovrapposizioni di pi\'u stati prodotto.
    \end{block}
\end{frame}

\begin{frame}{Cooper pairs}
    \begin{itemize}
        \item Model of pair of electrons added to Fermi sea at $T=0$, and they interact with each other but not with sea.
        \item Weak attraction can bind pairs of Fermions into bound state: Because of the Fermi sea, even an arbitrarily weak attraction can lower the energy of two fermions below $2E_F$, producing a bound Cooper pair.
        \item Two fermions must have equal and opposite momenta:
            \begin{columns}[T]
                \begin{column}{0.5\textwidth}
                    $\vec{k}$ momento relativo dei due fermioni, $\vec{K}$ momento totale della coppia
                    \begin{align*}
                        &\vec{k}_1=\vec{k}+\frac{\vec{K}}{1},\ \vec{k}_2=-\vec{k}+\frac{K}{2}
                    \end{align*}
                \end{column}
                \begin{column}{0.5\textwidth}
                    In BCS standard questo passaggio \'e giustificato da energetica e simmetria: si accoppiano stati time-reversed $(\vec{k},\shortuparrow)$, $(-\vec{k},\shortdownarrow)$.
                \end{column}
            \end{columns}
                    \begin{align*}
                        &E_{kin}=\frac{\hbar^2}{2m}[(\vec{k}+\frac{\vec{K}}{2})^2+(-\vec{k}+\frac{\vec{K}}{2})^2]=\frac{\hbar^2k^2}{m}+\frac{\vec{K}^2}{4m}\tag{Energia minima per $K=0$}
                    \end{align*}
                \item Funzione d'onda orbitale della forma $\psi_0(\vec{r}_1,\vec{r}_2)=\sum_{\vec{k}}\exp{i\scap{k}{r}_1}\exp{-i\scap{k}{r}_2}$
                \item Antisymmetry of total wavefunction to exchange of pair of fermions: searching for bound state we use $\cos{\vec{k}(\vec{r}_2-\vec{r}_2)}$ hence we have to use the antisimmetric singlet of spin as it gives higher probability to found the two particles near each other: $\psi_0(\vec{r}_1-\vec{r}_2)=[\sum_{k>k_F}g_{\vec{k}}\cos{\vec{k}(\vec{r}_1-\vec{r}_2)}](\alpha_1\beta_2-\beta_1\alpha_2)$
                \item Inserting into \schr{} eq: $V_{kk'}=\invers{\Omega}\int V(r)\exp{i(\vec{k'}-\vec{k})\cdot\vec{r}}$
                    \begin{columns}[T]
                        \begin{column}{0.5\textwidth}
                    \begin{equation*}
                        (E-2\epsilon_k)g_k=\sum_{k'>k_F}V_{kk'}g_{k'}
                    \end{equation*}
                        \end{column}
                        \begin{column}{0.5\textwidth}
                            \begin{align*}
                                &(H_0+V)\psi=E\psi,\ H_0=-\frac{\hbar^2}{2m}(\nabla_1^2+\nabla_2^2)\\
                                &H_0\phi_k=2\epsilon_k\phi_k
                            \end{align*}
                        \end{column}
                    \end{columns}
                    
                    con $\epsilon_{k}=\frac{\hbar^2k^2}{2m}$ energia onda piana imperturbata. Per $V=0$ autostato \'e coppia con $\vec{k}$ ed $E=2\epsilon_{k}$. Quando $V\neq0$ l'interazione collega $(\vec{k},-\vec{k})\leftrightarrow(\vec{k'},-\vec{k'})$
    \end{itemize}
\end{frame}

\begin{frame}{Bound pair state - Cooper Approx}
    \begin{itemize}
        \item If a set of $g_k$ exists satisfying $(E-2\epsilon_k)g_k=\sum_{k'>k_F}V_{kk'}g_{k'}$: a bound-state exists.
        \item Cooper Approx: $V_{kk'}=$\lbt{-V\quad E_F<E<E_F+\hbar\omega_c}{0}: due elettroni indipendenti posti appena sopra il mare di Fermi hanno energia $2E_F$ quindi binding energy $E_B=2E_F-E>0$.
        \item Summing both side of bound-state energy equation and using Copper Approx: $\frac{1}{V}=\sum_{k>k_F}\invers{[2\epsilon_k-E]}$.
        \item Passiamo al continuo di livelli energetici: $\sum\to\int N(\epsilon)\,d\epsilon$ con $N(0)$ densit\'a di stati attorno all'energia di Fermi per elettroni con 1 orientazione di spin: $\hbar\omega_c$ piccolo considero $N(\epsilon)\approx N(E_F)\approx N(0)$, quindi
            \begin{align*}
                &\frac{1}{V}=N(0)\int_{E_F}^{E_F+\hbar\omega_c}\frac{d\epsilon}{2\epsilon-E}=\frac{1}{2}N(0)\ln{\frac{2E_F-E+2\hbar\omega_c}{2E_F-E}}
            \end{align*}
        \item Weak coupling approx $VN(0)\ll1$ (in most classic superconductors it's found that $N(0)V<0.3$)
            \begin{align*}
                &E\approx 2E_F-2\hbar\omega_c\exp{-\frac{2}{N(0)V}}
            \end{align*}
            cio\'e esiste stato legato con energia negativa risp. superficie di Fermi costituito da elettroni con $k>k_F$.
        \item Electron states within range $\approx E'$ above $E_F$ are those most involved in bound states formation:
            \begin{columns}[T]
                \begin{column}{0.5\textwidth}
                    \begin{align*}
                        &\psi(r)\propto\sum_{k\geq k_F}\frac{\cos{\scap{k}{r}}}{2\xi_{\vec{k}}+E'},\ \xi_{\vec{k}}=\epsilon_{\vec{k}}-E_F,\ E'=2E_F-E>0\\
                        &E'\approx2\hbar\omega_c\exp{-\frac{2}{N(0)V}}\ll\hbar\omega_c
                    \end{align*}
                \end{column}
                \begin{column}{0.5\textwidth}
                    Spherical symmetry: S-state as well singlet spin state - Weight factor $\invers{2\xi_{k}+E'}$ has maximum at $\frac{1}{E'}$ where $\xi_k=0$. La regione vicina al cutoff $E_F+\hbar\omega_c$ \'e molto lontana dagli stati che dominano la coppia; superfluido neutronico non ha cutoff fononico dei superconduttori.
                \end{column}
            \end{columns}
        \item Size of bound pairs (Pippard): Coppia molto selettiva in $k$ \'e distante in $r$ - $\Delta r\geq \frac{1}{\Delta k}\approx\frac{\hbar \nu_F}{E'}$.
                    \begin{align*}
                        &\epsilon_k-E_F\approx\TDy{k}{\epsilon}|_{k_F}(k-k_F),\ \frac{1}{\hbar}\TDy{k}{\epsilon}=v_F:\ \Rightarrow \xi_k=\epsilon_k-E_F=\hbar v_F(k-k_F)\Rightarrow \Delta\xi\approx E'\Rightarrow \Delta k\approx \frac{E'}{\hbar v_f}
                    \end{align*}
    \end{itemize}
\end{frame}

\subsection{Superfluidity: BCS}

\subsection{Superfluidity: Landau-Ginzburg}

